\section{Proof of the identity}\label{sec:identity}

In this section, we prove the identity stated in
Theorem~\ref{thm:identity}. The proof proceeds by partitioning
the sum $D(x) = \sum_{n \leq x} d(n)$ according to the value of
$\omega(n)$, and then applying an inclusion--exclusion argument
to the composite part.

\begin{proof}[Proof of Theorem~\ref{thm:identity}]
Partition the integers $n$ with $1 \leq n \leq x$ into three
disjoint classes:
\begin{itemize}
    \item[(i)] $n = 1$ (equivalently, $\omega(n) = 0$);
    \item[(ii)] prime powers, i.e., $\omega(n) = 1$;
    \item[(iii)] composite integers with $\omega(n) \geq 2$.
\end{itemize}
Writing $D(x)$ as the sum of contributions from each class, we
have
\begin{equation}
    D(x) = d(1) + \sum_{\substack{n \leq x \\ \omega(n) = 1}} d(n)
               + \sum_{\substack{n \leq x \\ \omega(n) \geq 2}} d(n).
\end{equation}

Since $d(1) = 1$, the first term equals $1$.

For the second term, every integer $n$ with $\omega(n) = 1$ is of
the form $n = p^i$ for a unique prime $p$ and a unique integer
$i \geq 1$, and $d(p^i) = i + 1$. The number of prime powers
$p^i \leq x$ with exponent exactly $i$ equals $\pi(x^{1/i})$.
Hence
\begin{equation}
    \sum_{\substack{n \leq x \\ \omega(n) = 1}} d(n)
    = \sum_{i \geq 1} (i+1) \pi(x^{1/i}).
\end{equation}

For the third term, we apply inclusion--exclusion on the set of
distinct prime divisors of $n$. For each integer $j \geq 2$,
recall the definition
\[
    T_j(x) = \sum_{n \leq x} d(n) \binom{\omega(n)}{j}.
\]
We claim that
\begin{equation}
    \sum_{\substack{n \leq x \\ \omega(n) \geq 2}} d(n)
    = \sum_{j \geq 2} (-1)^j (j-1) T_j(x).
\end{equation}

To prove~(2.3), interchange the order of summation on the
right-hand side:
\[
    \sum_{j \geq 2} (-1)^j (j-1) T_j(x)
    = \sum_{n \leq x} d(n) \sum_{j \geq 2} (-1)^j (j-1)
      \binom{\omega(n)}{j}.
\]
For $n$ with $\omega(n) \in \{0, 1\}$, we have
$\binom{\omega(n)}{j} = 0$ for all $j \geq 2$, so such $n$ do not
contribute. For $n$ with $\omega(n) = m \geq 2$, we use the
following combinatorial identity.

\begin{lemma}
For every integer $m \geq 2$,
\begin{equation}
    \sum_{j=2}^{m} (-1)^j (j-1) \binom{m}{j} = 1.
\end{equation}
\end{lemma}

\begin{proof}[Proof of the lemma]
Starting from the standard identities
\[
    \sum_{j=0}^{m} (-1)^j \binom{m}{j} = 0
    \quad \text{and} \quad
    \sum_{j=0}^{m} (-1)^j \, j \binom{m}{j} =\left.\frac{d}{dx}(1-x)^m\right|_{x=1}=0 \ \text{(for } m \geq 2),
\]
we combine them to obtain
\[
    \sum_{j=0}^{m} (-1)^j (j-1) \binom{m}{j} = 0
    \quad \text{(for } m \geq 2).
\]
Separating the $j = 0$ and $j = 1$ terms:
\[
    1 \cdot (-1) + 0 \cdot \binom{m}{1}
    + \sum_{j=2}^{m} (-1)^j (j-1) \binom{m}{j} = 0,
\]
which yields
\[
    \sum_{j=2}^{m} (-1)^j (j-1) \binom{m}{j} = 1. \qedhere
\]
\end{proof}

Applying the lemma, we obtain
\[
    \sum_{n \leq x} d(n) \sum_{j \geq 2} (-1)^j (j-1)
    \binom{\omega(n)}{j}
    = \sum_{\substack{n \leq x \\ \omega(n) \geq 2}} d(n) \cdot 1,
\]
which establishes~(2.3).

Combining the contributions from all three classes yields
\[
    D(x) = 1 + \sum_{i \geq 1} (i+1) \pi(x^{1/i})
             + \sum_{j \geq 2} (-1)^j (j-1) T_j(x),
\]
completing the proof.
\end{proof}

